% Folder 156-8, batch 07 (archivists' pages 121-126).
% Pass: Opus 5.5 (claude-opus-5-5), 2026-10-02 — first pass, unchecked against the pages by a human.
%
% Grothendieck's own pagination runs 113-118 on these six leaves. The pencilled
% archive numbers on the leaves read 122-127, one ahead of the page numbering
% used here (the folder's yellow cover is counted by the pencil, not by us).
\documentclass[11pt,a4paper]{article}
\input{../preamble/grothendieck}

\folder{156-8}
\batch{7}
\pages{121}{126}
\foldertitle{[Chapitre] VIII. Analysis situs (quatrième mouture) : notes manuscrites (26/06-04/07/1986).}
\dating{1986}
\watermark{Édition de démonstration}

\begin{document}

\page{121}
\note{p. 113 de l'auteur. Le numéro d'archive porté au crayon sur ce feuillet est 122, et ainsi de suite jusqu'à 127 : il précède d'une unité la numérotation suivie ici, sur les six feuillets. La page s'ouvre au milieu d'une énumération de données commencée avant ce lot ; le « 1°) » n'en est pas ici.}

\note{Conventions de ce lot. La lettre $\Phi$ (la préfigure, et ses parties $\Phi_0$, $\Phi_1$, $\Phi_\Delta$, $\Phi'^*$\,…) est écrite partout avec un trait horizontal sous la lettre ; on transcrit simplement $\Phi$. Le signe d'ordre entre segments, un « $<$ » suivi d'un trait vertical, est rendu $\mathrel{<\!|}$. $\mathcal{M}$, $\mathcal{M}_1$ et $\mathcal{L}$ sont écrits en cursive.}

\struck{2°) Dans $\delta\Psi = \partial\Psi$,}

Les sommets \struck{isolés sont} intérieurs i.e. \struck{ceux de} \add{les} éléments de
\[
  \Delta_{\mathrm{int}} = (\Delta \setminus \partial\Psi) \cap \mathrm{Omb}(\Psi) = \Delta \cap \mathrm{Omb}^\circ(\Psi)
\]
\note{après le premier $\Delta$ de la parenthèse, un signe noirci, illisible.}

Les sommets isolés sont \struck{ceux de} \uncertain{les autres} sommets \struck{$\Delta \setminus \partial\Psi$} de $\Delta \setminus \partial\Psi$,
\[
  \Delta_{\mathrm{is}} = \text{\struck{$\Delta \setminus (\partial\Psi \cup \Delta$}}\; (\Delta \setminus \partial\Psi) \setminus \Delta_{\mathrm{int}}
\]

\struck{2°) \ill{}, d\ill{} d\ill{}, dans $\Delta_{\mathrm{int}}$,} \struck{une partie} \add{donc} une \uncertain{décomposition} en
\[
  \Delta = \partial\Psi \amalg \Delta_{\mathrm{int}} \amalg \Delta_{\mathrm{is}}
\]

2°) Se donner, dans $\Delta_{\mathrm{int}}$, les sommets qui appartiennent à $\Phi_0$ (ces sommets intérieurs propres, ou sommets de redondance ou de \ill{}), ou ce qui revient au même, les autres : les sommets lacunaires
\[
  \boxed{\Delta_{\mathrm{lac}}} \subset \Delta_{\mathrm{int}}
\]
de sorte que
\[
  \Delta_{\mathrm{int}} = \Delta_{\mathrm{lac}} \amalg \Delta_{\mathrm{red}} .
\]

3°) Dans $\partial\Psi$ \add{$(= \delta\Psi)$}, dire quels sommets doivent être dans $\Phi_0$ \struck{et lesquels} (ces sommets bords propres de $\Phi$) et lesquels non, i.e. se donner une partie
\[
  \partial_0 \subset \partial\Psi .
\]
\note{au-dessus de $\partial_0$, un premier symbole biffé.}

En \uncertain{résumé}
\note{la phrase s'arrête là ; la page suivante repart sur une « Proposition ».}

\page{122}
\note{p. 114 de l'auteur.}

\textbf{Proposition.} Une préfigure quelconque est déterminée par les données suivantes :

1°) Une figure $\Psi$ \add{sans sommets isolés et, à composantes connexes} \add{\struck{(i.e. $\delta\Psi = \partial\Phi$} i.e. telle que $\delta\Psi = \partial\Phi$)} des \struck{multi} segments, \add{déterminée par}
\[
  D = (X_1 \mathrel{<\!|} X_2 \mathrel{<\!|} \text{\struck{$X_3$}} \cdots X_n) \in \mathrm{Drap}^*(\mathcal{M}_1) .
\]

\struck{2°) Une} \struck{tel} par
\[
  \Psi = \Phi_D = \bigcup_{X \in D} \widetilde{X}_{\xi} = \bigcup_{1 \leq i \leq n} \widetilde{X}_i
\]
\note{l'indice de $\widetilde{X}$ dans la première réunion ressemble à un $\xi$ ; lecture douteuse.}

(réunion disjointe). On a \uncertain{aussi}
\[
  \partial\Psi = \delta\Psi = \bigcup_{X \in D} \partial X = \bigcup_{1 \leq i \leq n} \partial X_i
\]
(réunion disjointe), soit
\[
  \delta\Psi = \{a_1, b_1, a_2, b_2, \ldots, a_n, b_n\}
\]
si
\[
  \partial X_i = \{a_i, b_i\} .
\]

On peut dire aussi que $\Psi$ est déterminée par la suite \struck{strictement} croissante de cardinal pair
\[
  a_1 < b_1 < a_2 < b_2 \cdots < a_n < b_n ,
\]
et on aura
\[
  X_i = S_{a_i, b_i} \qquad 1 \leq i \leq n .
\]

\marginal{2°) ici, cf. \uncertain{p.~115} \ill{}}
\note{renvoi en marge gauche, au-dessus d'un trait de séparation : le 2°) de la Proposition est rédigé p. 115 de l'auteur (page 123 de la présente numérotation).}

3°) Un drapeau $\Delta \in \mathrm{Drap}^*(\mathcal{L})$,
\[
  \Delta \supset \partial\Psi
\]

\struck{3°) Une partie $\Delta_{\mathrm{lac}}$}

d'où on déduit une partie
\note{la phrase se poursuit en tête de la page suivante.}

\page{123}
\note{p. 115 de l'auteur.}

\struck{$\Delta_{\mathrm{is}} =$ \ill{} $\Delta \setminus \Delta \cap \mathrm{Omb}^\circ(\Psi_1)$}
\note{ligne biffée ; sous $\Delta \cap \mathrm{Omb}^\circ(\Psi_1)$, une accolade marquée $\Delta_{\mathrm{int}}$.}

Posant
\[
  \Delta_{\mathrm{int}} = \Delta \cap \mathrm{Omb}^\circ(\Psi_1) \subset \Delta \setminus \partial\Psi
\]
(\uncertain{corr.} aux pts de $\Delta$ \emph{intérieurs} \struck{\ill{}} \add{\ill{}} des intervalles $S_{a_i, b_i}$) on aura
\[
  \Delta = \Delta_{\mathrm{is}} \amalg \Delta_{\mathrm{int}} \amalg \partial\Psi
\]
où
\[
  \Delta_{\mathrm{is}} \overset{\mathrm{def}}{=} (\Delta \setminus \partial\Psi) \setminus \Delta_{\mathrm{int}}
\]

\struck{3°)} \add{4°)} Une partie
\[
  \Delta_{\mathrm{lac}} \subset \Delta_{\mathrm{int}}
\]
\marginal{\struck{\ill{}}}
\note{en marge gauche, un mot biffé muni d'une flèche descendante.}

et on \struck{pose} aura donc
\[
  \Delta_{\mathrm{int}} = \Delta_{\mathrm{lac}} \amalg \Delta_{\mathrm{red}}
\]
où
\[
  \Delta_{\mathrm{red}} \overset{\mathrm{def}}{=} \Delta_{\mathrm{int}} \setminus \Delta_{\mathrm{lac}} .
\]

\struck{3°)} \add{2°)} Une partie
\[
  \partial_0 \subset \partial\Psi
\]
de $\partial\Psi$.

À ces données, on associe la préfigure $\Phi$ définie ainsi :
\[
  \begin{cases}
    \Phi_0 = \underbrace{\partial_0}_{\text{sommets bord propres}} \cup \underbrace{\Delta_{\mathrm{red}}}_{\text{sommets redondants}} \cup \Delta_{\mathrm{is}} \\[1ex]
    \Phi_1 = \{ X \in \underbrace{\Phi_\Delta \cap \mathcal{M}_1}_{\text{intervalles interstitiels de } \Delta} \mid X \in \mathrm{Omb}(\Psi) \}
  \end{cases}
\]
\note{devant $\Phi_0$, un signe biffé.}

\page{124}
\note{p. 116 de l'auteur.}

\struck{tel.}

et on aura
\[
  \begin{gathered}
    \delta\Phi = \Delta \\
    \partial\Phi = \partial\Psi , \quad \partial_0\Phi = \partial_0 , \quad \partial_{\mathrm{imp}}\Phi = \partial\Psi \setminus \partial_0 \\
    \Phi_{\mathrm{is}} = \Delta_{\mathrm{is}} \\
    \Phi_{\mathrm{red}} = \Delta_{\mathrm{red}} \\
    \Phi_{\mathrm{lac}} = \Delta_{\mathrm{lac}}
  \end{gathered}
\]

Toute \uncertain{préfigure} s'obtient \add{de façon unique} par des données 1°) 2°) 3°) 4°) \struck{(i.e. $\Psi$,} \struck{$\Delta_{\mathrm{is}}$} \struck{$\partial_0 \subset \partial\Psi \subset \Delta$}
\[
  \Psi \quad \text{ou} \quad X_1 \mathrel{<\!|} X_2 \mathrel{<\!|} X_3 \cdots \mathrel{<\!|} X_n
\]
\[
  \text{ou} \quad a_1 < b_1 < a_2 < b_2 \cdots < a_n < b_n
\]
et
\[
  \boxed{\partial_0} \subset \partial\Psi \subset \boxed{\Delta} \supset \Delta_{\mathrm{lac}}
\]
\note{$\Delta_{\mathrm{lac}}$ est écrit sous $\Delta$, relié par un signe $\cup$ vertical ; $\partial_0$ et $\Delta$ sont entourés.}

\struck{$\Delta_{\mathrm{int}} \subset \Delta_{\mathrm{lac}} \cap \mathrm{Omb}^\circ(\Psi_1)$} \struck{$= (\Delta \setminus \partial\Psi) \cap \mathrm{Omb}(\Psi_1)$}
\note{ligne biffée, la partie gauche raturée en hachures ; lecture du début douteuse.}

avec la condition que
\[
  \Delta \in \mathrm{Drap}^*(\mathcal{L}) ,
\]
et
\[
  \Delta_{\mathrm{lac}} \subset \mathrm{Omb}^\circ(\Psi_1)\, \text{\struck{\ill{}}} \cap \Delta
  \quad \bigl(= \mathrm{Omb}(\Psi_1) \cap (\Delta \setminus \partial\Psi)\bigr) .
\]

La préfigure $\Phi$ est irredondante sss
\[
  \Delta_{\mathrm{red}} = \emptyset
\]
i.e. sss \struck{$\Delta_{\mathrm{lac}}$ \ill{}}
\[
  \Delta_{\mathrm{lac}} = \Delta \cap \mathrm{Omb}^\circ(\Psi_1) ,
\]
(où $\mathrm{Omb}^\circ(\Psi) = \bigcup_{1 \leq i \leq n} \mathrm{Omb}^\circ(X_i)$ (réunion disjointe)).

\textbf{Corollaire 1.} Se donner une \add{préfigure} $\Phi$ irredondante revient au même que se donner $\Psi$, et $\partial_0$, $\Delta$ avec \struck{\ill{}}
\[
  \partial_0 \subset \partial\Psi \subset \Delta
\]
\ill{} $\Delta \in \mathrm{Drap}^*(\mathcal{L})$.
\note{l'énoncé du corollaire est marqué d'un trait vertical en marge gauche.}

\page{125}
\note{p. 117 de l'auteur.}

\textbf{Corollaire 2} (Pour mémoire) \add{si $\Phi$ irredondante} Moyennant Hyp. \ill{}
\[
  \mathrm{Supp}^\circ(\Phi) = \mathrm{Omb}^\circ(\Phi_1) \cup \text{\struck{\ill{}\,$\Delta_{\mathrm{is}}$}}\; \Phi_0
\]
\note{au-dessus du membre de droite, inséré : « $\mathrm{Omb}^\circ(\Phi) =$ ». La ligne suivante, entièrement biffée, commençait par « $=$ » et finissait par « $\partial_0 \cup \Delta_{\mathrm{is}}$ », avec un indice $1 \leq i \leq n$ ; elle n'est pas lisible.}

(réunions disjointes).

\textbf{Corollaire 3.} Considérons la figure $\Phi'^*$ \struck{formée} \add{définie}
\struck{1°) des \struck{figures} \add{intervalles} interstitiels de $\Delta$ qui \ill{} pas dans $\Phi$ i.e. qui \ill{} dans $\mathrm{Omb}(\Psi)$, et $\Phi'^*_0 = \delta\Delta \setminus \Phi_0 = \partial'_0$}
\note{ce passage est annulé par un contour fermé et une série de traits obliques ; il est repris en dessous.}

\[
  \Phi'^*_0 = \underbrace{\delta\Phi}_{\Delta} \setminus \Phi_0 = \Delta \setminus \Phi_0 = \Delta_{\mathrm{lac}} \amalg \partial'_0
\]
où je pose
\[
  \partial'_0 \overset{\mathrm{def}}{=} \partial\Phi \setminus \partial_0\Phi = \partial\Psi \setminus \partial_0
\]
et
\[
  \Phi'^*_1 = \text{\struck{$\{ X \in \Phi_\Delta \cap \mathcal{M}_1 \mid X \notin \Phi_1$}\,\struck{\ill{}}}\; \}
\]
\note{sous la formule, un signe biffé, puis un trait ondulé de séparation.}

\marginal{Rappel :
$\Delta = \Delta_{\mathrm{is}} \amalg \Delta_{\mathrm{int}} \amalg \partial\Psi$,
$\Phi_0 = \Delta_{\mathrm{is}} \amalg \Delta_{\mathrm{red}} \amalg \partial_0$,
où $\Delta_{\mathrm{int}} = \Delta_{\mathrm{red}} \amalg \Delta_{\mathrm{lac}}$,
$\partial\Psi = \partial_0 \amalg \partial'_0$.}

Ça a l'air franchement compliqué comme présentation. Ne vaut-il pas mieux dire ceci. Se donner une préfigure $\Phi$, c'est se donner
\begin{enumerate}
  \item[1°)] $\Delta = \delta\Phi \in \mathrm{Drap}^*(\mathcal{L})$
  \item[2°)] Une partie $\Phi_1 \subset (\Phi_\Delta)_1 =$ ens. des intervalles interstitiels de $\Delta$
  \item[3°)] Une partie $\Phi_0 \subset (\Phi_\Delta)_0 = \Delta$
\end{enumerate}

\page{126}
\note{p. 118 de l'auteur.}

\uncertain{Err}. Pour de telles données, la réunion $\Phi_0 \cup \Phi_1$ est bien une figure, mais $\delta(\Phi_0 \cup \Phi_1) \subset \Delta$ et l'inclusion peut être stricte, (exemple : $\Phi_0 = \Phi_1 = \emptyset$ !) et il faut donc de plus poser une condition qui implique que $\delta(\Phi_0 \cup \Phi_1) = \Delta$), donc \uncertain{assurer} la condition
\[
  \text{\uncertain{Pur}\ill{}}_\Delta : \quad \forall s \in \Delta \setminus \Phi_0 ,\ \exists X \in \Phi_1 \text{ t.q. } s \lhd X ,
\]
en d'autres termes
\[
  \Delta = \Phi_0 \cup \bigcup_{X \in \Phi_1} \partial X .
\]
\note{le nom de la condition n'est pas sûr.}

\marginal{Si $\Delta = \emptyset$, $\Phi_\Delta = \emptyset$, $\Phi = \emptyset$ \uncertain{trivial}. Et \ill{}, on suppose $\Delta \neq \emptyset$}

Considérons alors
\[
  \Phi'^*_0 = \text{\struck{$\Delta \setminus \Phi$}}\; (\Phi_\Delta)_0 \setminus \Phi_0 = \Delta \setminus \Phi_0
\]
\[
  \Phi'^*_1 = (\Phi_\Delta)_1 \setminus \Phi_1
\]

Soit de plus, \struck{si $\Delta \neq \emptyset$ (i.e. $\Phi \neq \emptyset$)},
\[
  a = \mathrm{or}(\Delta) , \quad b = \mathrm{ex}(\Delta)
\]
et soit
\[
  \begin{gathered}
    \mathcal{M}_{<a} = \{ X \in \mathcal{M} \mid X < a \} \\
    \mathcal{M}_{>b} = \{ X \in \mathcal{M} \mid X > b \}
  \end{gathered}
\]

Par Hyp., $\mathcal{M}_{<a}$ est $\emptyset$ sss $a$ est un plus petit élément de $\mathcal{L}$, et $\mathcal{M}_{>b}$ est vide sss $b$ est un plus grand élément de $\mathcal{L}$. Ceci dit, on trouve aussitôt
\[
  \mathrm{Cosupp}^\circ(\Phi) = \mathrm{Omb}^\circ(\Phi'^*_1) \cup \mathcal{M}_{<a} \cup \mathcal{M}_{>b}
\]
(réunion disjointe) tandis que
\[
  \mathrm{Omb}^\circ(\Phi'^*) \overset{\mathrm{def}}{=} \bigcup_{X \in \Phi'^*} \mathrm{Omb}^\circ(X) \qquad \text{réunion disjointe}
\]
\note{l'indice $1$ de $\Phi'^*_1$ dans $\mathrm{Cosupp}^\circ$ est récrit sur un autre chiffre. Le lot s'arrête ici ; la comparaison entre $\mathrm{Cosupp}^\circ(\Phi)$ et $\mathrm{Omb}^\circ(\Phi'^*)$ se poursuit au-delà.}

\end{document}
